\section{The endpoint and the response of the effective coefficient}
\label{ep:section}

Throughout this section $J$ is fixed, $v=v_J$, $A_0=-\Delta_J/v^2$,
and $k=\beta/v^2$.  Write $a(k)$ for the coefficient constructed in
Sections~\ref{st:section} and~\ref{an:section}, and put
$k_c=\beta_c(J)/v^2$.  Corollary~\ref{an:physical-criterion},
Proposition~\ref{an:coefficient-regularity}, and
Corollary~\ref{wh:high-temperature} give
\[
 0<k_c<\infty,\qquad
 \{k:a(k)>0\}=[k_c,\infty),\qquad
 a(k)\in\{0\}\cup[8\pi,\infty).
\]
On its positive set $a$ is strictly increasing and right continuous.
Right continuity here follows also directly by combining monotonicity
with upper semicontinuity.  We shall prove $a(k_c)=8\pi$,
differentiability on $(k_c,\infty)$, and an infinite right derivative
and right difference quotient at $k_c$.

A reference coefficient, denoted throughout by $a'$, is held fixed
when differentiating in $k$.  In particular a prime on $a'$ is part
of its name, not a derivative.  We write $\frac{d}{dk}a(k)$ for the
derivative of the physical coefficient.  Set $\alpha=\sqrt{a'}$ and
$\lambda=L^{2-a'/(4\pi)}$.

The critical small-activity Coulomb-gas trajectory constructed by
Falco \citep{Falco2012} is a rigorous predecessor of the marginal
flow analysis.  Here the preceding entry theorem and the physical
terminal diagnostics identify which trajectory is selected by the
original fixed-range height model.  The finite-volume response then
connects that trajectory to the endpoint slope.

\subsection{Physical terminal diagnostics}

Fix $k_h$ with $a_h=a(k_h)>0$.  Choose the downstream norm parameters
in Proposition~\ref{wh:map} and Theorem~\ref{en:entry} first.  The entry
construction then chooses $R$, $P=R^5$, a finite large-field threshold,
and $s$, in this order, before the parameter neighborhood, and starts
the long maps at $m=qRs$.  We use the observation
$Y=B_s h+\eta$ of Theorem~\ref{en:entry}, where the noises have
variance one.  All choices are fixed on a sufficiently small
neighborhood of $(k_h,a_h)$.  Near $8\pi$ retain the gradient and
fundamental coordinates $(t,z)$ on the allowed one-sided interval
$8\pi\le a'\le8\pi+\delta_{\rm ref}$; on a compact reference interval
strictly above $8\pi$ choose $L$ for its gap and retain only $t$.
All these continuations use the fixed-$J$ entry shells
$\Gamma_u^{\rm en}$ and their exact multipliers
$\lambda_j^{\rm en}$ from Lemma~\ref{wh:kernels}.

Choose a sufficiently large fixed dyadic integer $D$ and use actual
tori of sides
\begin{equation}\label{ep:stop-sizes}
 n_j=Dm_j,\qquad m_j=L^jm .
\end{equation}
The choice of $D$ exceeds the causal neighborhoods in
Proposition~\ref{wh:map}.  Consequently all singleton coefficients
through the stop at $m_j$ agree exactly with the plane coefficients.
These sides satisfy the entry divisibility condition, since
$n_j/s=DqR L^j$ is a multiple of $qR$ with the chosen grid origin.
The torus remainder need not agree with the plane remainder:
noncopying supports are estimated by the same contraction.
Winding links keep their full covers, as in Theorem~\ref{en:entry},
so they cannot change a singleton in a padded unwrapped
neighborhood.  When $n_j<sP$, only the plane steps are used.
Their number is at most
$J_R=\max\{0,\lceil\log_L(R^4/(Dq))\rceil\}$, independently of $s$.
Equation~\eqref{en:placement} gives an entry norm at most $Ce^{-cR}$,
so the bounded plane map amplifies it by at most
$C_0^{J_R}=R^{O(1)}$; increasing $R$ keeps all these steps small.
For each valid stopping size $n_j\ge sP$, its own exact torus
history starts at $m$ and runs through all steps up to $j$.
No identity on a smaller invalid torus is needed.

Let $f$ be the real first horizontal cosine on the unit torus and
let $f_n$ be its cell-center samples, weighted by the cell area.
Put
\[
 X_n=(Y,f_n),\qquad
 g=(f,(-\Delta_{\T})^{-1}f)>0 .
\]
The height component is $(h,B_s^*f_n)$.  Its difference from the
ordinary sampled smooth test has bare inverse-energy norm tending
to zero as $s/n\to0$.  Also
$\Var(\eta,f_n)=O((s/n)^2)$.  The field and Laplace conclusions
of Theorem~\ref{an:dirichlet-torus} therefore apply to $X_n$,
including convergence of every fixed moment.

For a real source $w$ near zero, Gaussian completion translates the
observation argument of the entire interaction by its reference
covariance applied to $wf_n$.  This translate equals $B_s\alpha F$
for a microscopic cosine $F$: translation and reflection symmetry
make the translated cell function a cosine, and the average of a
microscopic sampled cosine is that cell cosine times a nonzero
scalar.  The scalar is uniformly bounded at the sizes
\eqref{ep:stop-sizes}.  Its energy satisfies
\begin{equation}\label{ep:F-energy}
 (F,A_0F)=a'w^2g+o(1)
 \quad(s/n\to0).
\end{equation}
All its scaled derivatives needed in the point norm are bounded.
Field-independent extracted constants cancel in the ratio.

Write $\mathcal R_j^{\rm tor}$ for the stopped torus remainder and
$\nu_j=|t_j|+|z_j|+\|\mathcal R_j^{\rm tor}\|$, omitting $z$
in the one-coordinate scheme.  Terminal integration has only
$D^2$ blocks, and its regulators integrate with reserve by
Lemma~\ref{en:regulator}.  Expanding the exact terminal
partition ratio at zero activity gives
\begin{equation}\label{ep:gradient-diagnostic}
 \log\E e^{wX_{n_j}}-\log\E_{a'}e^{wX_{n_j}}
 =\frac12t_j(F,A_0F)
       +O\bigl(\|\mathcal R_j^{\rm tor}\|+\nu_j^2\bigr).
\end{equation}
The bound is uniform for $w$ in a fixed complex disk; thus it also
holds after taking any fixed number of source derivatives.  The
constant mean integral kills the linear fundamental term.

A second diagnostic detects that term.  Write
$\chi_n=e^{i\bar Y}$, where $\bar Y$ is the spatial mean modulo
$2\pi$.  Let $\zeta_{{\rm high},j}$ be the sum of the polynomial
Gaussian pieces integrated through side $m_j$, including their
constant-mode continuations, and define the variance in phase units by
\[
 \tau_{\rm av}
 =\Var\bigl(\alpha\,\overline{B_s\zeta_{{\rm high},j}}+\bar\eta\bigr).
\]
Here $\bar\eta$ is the spatial mean of the observation noise.
This high mean is independent of all mean-free Gaussian variables
and of their source tilt.  Its variance is bounded: a polynomial
piece with upper side $u\le m_j$ contributes $O(u^2/n_j^2)$, the
sum over such sides is geometric, and
$\Var\bar\eta=(s/n_j)^2$.

Let $\mathcal Z_j$ be the stopped subset expansion with its common
field-independent factors removed.  Let $C$ be uniform on
$[0,2\pi/\alpha)$, let $\zeta_{*,j}$ have the remaining mean-free
terminal covariance, and let $F_j(w)$ be the field translate above
at side $n_j$.  Put
\[
 \begin{split}
 c_{n_j,a'}&=(f_{n_j},[I+a'B_sA_0^+B_s^*]f_{n_j}),\\
 \mathcal I_j(w)&=\E_{C,\zeta_{*,j}}
       \mathcal Z_j(C+\zeta_{*,j}+F_j(w)),\\
 \mathcal N_j(w)&=\E_{C,\zeta_{*,j}}
       \bigl[e^{i\alpha C}\mathcal Z_j(C+\zeta_{*,j}+F_j(w))\bigr].
 \end{split}
\]
Gaussian completion and the uniform mean integral give the exact
normalized mixed identity
\[
 e^{-c_{n_j,a'}w^2/2}\E[\chi_{n_j}e^{wX_{n_j}}]
       -\E\chi_{n_j}
 =e^{\tau_{\rm av}/2}
       \frac{\mathcal N_j(w)-\mathcal N_j(0)}{\mathcal I_j(0)} .
\]
Indeed absorbing the integrated high mean into $C$ leaves the
physical character unchanged, whereas inserting $e^{i\alpha C}$
before that change contributes the Gaussian characteristic
$e^{-\tau_{\rm av}/2}$.  This proves the positive sign in the
displayed prefactor.  The source annihilates constants, so it does
not change that argument.  The denominator is the common zero-source
normalizer; also $c_{n_j,a'}\to a'g$ as $s/n_j\to0$.

At zero activity $\mathcal I_j(0)=1$, and on the small ball it is
$1+O(\nu_j)$.  In $\mathcal N_j(w)-\mathcal N_j(0)$ the constant
and linear gradient terms vanish under the mean character.  Exact
integration of the linear cosine term therefore makes the preceding
mixed functional
$z_jM_j(w)+O(\|\mathcal R_j^{\rm tor}\|+\nu_j^2)$, where
\begin{equation}\label{ep:phase-multiplier}
 M_j(w)=\frac{D^2}{2}
 e^{\tau_{\rm av}/2-\alpha^2 D_{{\rm tail},j}/2}
 \operatorname{ave}_x\bigl(e^{-i\alpha F_j(w;x)}-1\bigr).
\end{equation}
Here $D_{{\rm tail},j}$ is the diagonal variance of the remaining
mean-free Gaussian field.  The entry terminal bound gives
$D_{{\rm tail},j}\le C_J(1+\log D)$, so it is uniformly bounded.
For a fixed
small nonzero real $w$, the average in
\eqref{ep:phase-multiplier} tends to a strictly negative number:
the spatial cosine has zero mean and a strictly positive second
moment.  In particular $|M_j(w)|\ge c_w>0$ for all sufficiently
small $s/n_j$.  The error in this mixed diagnostic is
again $O(\|\mathcal R_j^{\rm tor}\|+\nu_j^2)$.
These estimates, including their input derivatives at fixed $a'$,
follow by the analytic terminal expansion, rather than by
differentiating a limiting measure.

\begin{lemma}[Phase factorization at a true reference]
\label{ep:phase-factorization}
At $(k,a')=(k_h,a_h)$,
\[
 \E\bigl[\chi_n e^{wX_n}\bigr]
 -\E[\chi_n]\exp(a_hgw^2/2)\longrightarrow0
\]
locally uniformly for complex $w$.  The conclusion does not assert
that $\E\chi_n$ has a nonzero limit or that the torus mean is
asymptotically uniform.
\end{lemma}

\begin{proof}
We give the replica argument, since Gaussian convergence of $X_n$
alone would not imply this assertion.  First remove the independent
observation noise; its nonconstant test variance tends to zero,
and its mean characteristic is a bounded deterministic multiplier.
Take two independent pinned lattice fields $h^{(1)},h^{(2)}$.
The sum and difference are conditionally independent, with the
same centered law, when their common coordinatewise parity class
is fixed.  This is the elementary change of variables
$(h^{(1)},h^{(2)})\mapsto(h^{(1)}+h^{(2)},h^{(1)}-h^{(2)})$:
the quadratic energy splits and the two fields belong to the same
coset of twice the height lattice.

Let $V$ be the conditional second moment of the cosine test of the
sum field.  Expansion in the independent original copies gives
\[
 \E V=2\Var X_n,\qquad
 \Var V=2\bigl(\E X_n^4-3(\E X_n^2)^2\bigr).
\]
The last expression tends to zero by Gaussian fourth-moment
convergence.  Multiplication by a bounded character of the
difference field therefore replaces $V$ by its mean with error
tending to zero.  Define
\[
 F_n(t)=\E e^{i\bar h+itX_n}.
\]
The corresponding product expectation of the squared sum test and
the difference character is
$-2(F_n F_n''-(F_n')^2)$.  More explicitly, for complex $t$,
\[
 -2(F_nF_n''-(F_n')^2)
 =2\Var(X_n)F_n^2+
 \E[(V-2\Var(X_n))
      e^{i\bar D+itX_n(D)}].
\]
If $\E e^{tX_n}\le e^{B_0t^2/2}$ for real $t$, the
last expectation has modulus at most
$\sqrt{2\kappa_4(X_n)}e^{2B_0(\operatorname{Im}t)^2}$
by Cauchy--Schwarz.  Consequently every locally uniform
subsequential limit $F$ satisfies
\begin{equation}\label{ep:replica-ode}
 FF''-(F')^2=-a_hg F^2 .
\end{equation}
Bare domination supplies locally uniform bounds on the entire
functions and all their fixed derivatives.  The real transforms
are nonnegative by the Gaussian-lattice Fourier representation.
If $F$ is not identically zero, solve
$(\log F)''=-a_hg$ on any positive interval.  Analytic continuation
then makes $F$ an exponential quadratic everywhere; in particular
it has no isolated zero.  Reflection or translation by half a
cosine wavelength changes $X_n$ to $-X_n$ without changing
$\bar h$, so $F'(0)=0$.  Thus
$F(t)=F(0)e^{-a_hgt^2/2}$.  The identically zero limit has the same
factorization.  Analytic continuation to $t=-iw$ gives the stated
Laplace version.  Every subsequence has only limits of this form,
which proves the assertion for the original sequence.
\end{proof}

\begin{proposition}[No exit and vanishing at true references]
\label{ep:no-exit}
For every sufficiently small prescribed tube radius $\rho>0$,
the entry choices can be made so that the true-reference plane
trajectory stays in this tube at every step and tends to zero.
The stopped torus remainders also tend to zero.  If
$a(k_c)=8\pi$, the same entry choices keep all true-reference
trajectories $(k,a(k))$, $k\in[k_c,k_c+\delta)$, in that tube for
some $\delta>0$.  Each of these trajectories tends to zero;
uniform convergence in $k$ is not asserted here.
\end{proposition}

\begin{proof}
At the true reference both diagnostics tend to zero:
Equation~\eqref{ep:gradient-diagnostic} uses the structural
Laplace limit, and the mixed diagnostic uses
Lemma~\ref{ep:phase-factorization}.  Their discrepancies can be
made uniformly small for $n\ge sP$ by choosing $P$ large first and then $s$ sufficiently
large.  The block-profile error and the noise variance are
$O((s/n)^2)\le O(P^{-2})$, which explains this order.  Before the first candidate exit, the separate torus
recursion gives
\[
 \|\mathcal R_j^{\rm tor}\|
 \le C\theta^j\eps+C\rho^2
\]
with the same earlier plane singleton inputs.  At the candidate
exit the new activities are still in a fixed multiple of the
tube by the bounded map estimates.  The nonzero two diagnostic
multipliers therefore give
$|t_j|+|z_j|\le C(\delta_{\rm diag}+\eps+\rho^2)$.
Choose $\rho$ first, with $C\rho\ll1$, and then
$\delta_{\rm diag}+\eps\ll\rho$.  This excludes the exit.
The initial plane steps corresponding to sizes below $sP$ were
controlled just after Equation~\eqref{ep:stop-sizes}.

For convergence let $q$ be the limsup of the plane activity size.
Contraction bounds the limsup of its remainder, and separately
that of the stopped torus remainder, by $Cq^2$.  Taking limsups
in the two diagnostics gives $q\le C'q^2$.  The tube is chosen
smaller than $1/C'$, so $q=0$.

For the final assertion, right continuity gives
$a(k)\downarrow8\pi$ as $k\downarrow k_c$.  Fix a slightly
larger input $k_1>k_c$.  Finite-volume Gaussian comparisons
sandwich the mean-free Laplace functions for
$k_c\le k\le k_1$ between those at the two endpoints.
Their limiting Gaussian coefficients approach the same value
as $k_1\downarrow k_c$.  Entire-function bounds turn this into
uniform smallness of the requisite second and fourth moment
discrepancies.  The replica calculation in
Lemma~\ref{ep:phase-factorization} then gives the mixed
discrepancy uniformly as well: otherwise a violating sequence
has an entire subsequential limit obeying
\eqref{ep:replica-ode}, which is impossible.
Choose the auxiliary bracket endpoint $k_1$ sufficiently close to
$k_c$ first, then a common large-scale cutoff for the two fixed
endpoint laws and entry scales with $sP$ above that cutoff.
Finally shrink the actual input neighborhood to
$[k_c,k_c+\delta]\subseteq[k_c,k_1]$ as required by entry.
The preceding first-exit argument now holds throughout this
right neighborhood.  For each fixed
input its limiting discrepancy is zero, so the individual
limsup argument still proves convergence.
\end{proof}

\subsection{Exact Gaussian curves and positive response}

\begin{lemma}[Removing all pure gradient drift]
\label{ep:gaussian-curves}
Fix $J$, an allowed compact interval of reference coefficients, and
admissible long-map parameters, including $L,A,h_0,\kappa$, the
regulator margins $h_*,p_1,p_2$ when $e^V$ is used, and the terminal
clearance $D$.  There is a sufficiently large starting scale $m_{\min}$
and a disk $|d|<d_0$ such that, at every entry scale $m\ge m_{\min}$
and every reference in that interval, the long maps have exact Gaussian
trajectories satisfying
\[
 T_j^g(d)=d+O(d^2),\qquad Z_j^g(d)=0,\qquad
 \mathcal R_j^g(d)=O(d^2).
\]
The bounds are uniform in the step, the entry scale, and the reference,
and analytic in $d$.  The maps $T_j^g$ have analytic inverses on a
common smaller disk.  In particular both disks can be chosen before
the physical entry scale and its observation-cell size are selected.
In the coordinates
\begin{equation}\label{ep:normal-coordinates}
 u_j=(T_j^g)^{-1}(T_j),\qquad
 B_j=\mathcal R_j-\mathcal R_j^g(u_j),
\end{equation}
the exact pure Gaussian trajectory is $(u,0,0)$ at every step.
Its first coordinate is unchanged, so there is no nonlinear
forcing depending only on $u$.
\end{lemma}

\begin{proof}
At side $m$ initialize with the exponential of
$d\sum_B e_B^0/2$, with no height-lattice factor.  Expand it
by the same occupied-block algebra.  The singleton linear term
is $d e_B^0/2$ and the remainder is $O(d^2)$ in any prescribed
fixed activity weight, since the small exponential energy is
absorbed by the regulator reserve.  We check that this assertion is
uniform in the starting scale.  A direction of scaled point norm one
has $e_B^0(f)\le C_J$: there are $O(m^2)$ sites in a block and each
fixed-range difference is $O_J(m^{-1})$.  Polarization gives
$|De_B^0(\phi)[f]|\le C_J\sqrt{U_m}$ and
$|D^2e_B^0[f,g]|\le C_J$; all higher derivatives vanish.
The analytic derivative sum of a sufficiently small energy exponential
is therefore absorbed by a fixed part of $W_m^\kappa$, uniformly for
$m\ge m_{\min}$.  The occupied-block expansion with the fixed weight
$A$ has bounds $(C_A|d|)^{|X|}$ and, after removing its linear
singleton, $C_A|d|^2$.  The optional factor $e^V\ge1$ only improves
this initialization.  Common-shift invariance keeps every charged
coordinate zero.

Direct Gaussian completion bounds all terminal logarithmic
ratios under the smooth shift $F$ by $C|d|$, uniformly in
the stopping scale.  Indeed the tail covariance contracts
the kinetic energy, and the relevant inverse
$(I-dC^{1/2}A_0C^{1/2})^{-1}$ is uniformly bounded on a fixed
complex disk; determinants cancel from the ratio.
The kinetic contraction is the form inequality
$0\le C^{1/2}A_0C^{1/2}\le I$ for the remaining covariance, so
its resolvent norm is at most $(1-|d|)^{-1}$.  The selected kinetic
energy has shell and terminal expectation at most $C|X|$, uniformly
in the starting scale, by the positive-order kernel estimates and
Equation~\eqref{en:energy-contraction}.  Thus the terminal analytic
estimates use the same regulator reserve at every such scale.

Take diagnostic tori of side $n=DmL^j$ with the fixed $D$ chosen
large enough for locality.  The first Fourier mode used in the gradient
diagnostic has energy bounded above and below uniformly in $m,j$.
Indeed its symbol converges to the continuum symbol for
$m\ge m_{\min}$, and its $s$-cell averaging multiplier is bounded
away from zero when $s/n\le1/D$ (here $s\le m$).
The independent observation-noise contribution is uniformly bounded
as well.  No lower restriction $n\ge sP$ is needed for this Gaussian
initialization, because there are no entry sewing operations.
The gradient diagnostic consequently gives, before a first exit from
a fixed small tube, a bound $C|d|+C\|\mathcal R_j^{\rm tor}\|+C\rho^2$
on its scalar coordinate, with constants independent of $m,s,j$.
Contraction bounds the remainder by $C|d|^2+C\rho^2$ there.
Choose the tube radius first and then the common $d_0$ sufficiently
small.  The same first-exit argument works for complex $d$, using
absolute values of the analytic diagnostic, and excludes an exit
on that disk.  More quantitatively, on each parameter disk
$|d|\le\delta\le d_0$ take the tube radius $\rho=M\delta$, with
$M$ fixed larger than the diagnostic constant.  For sufficiently
small common $d_0$ the same estimates exclude exit from that tube.
This gives a uniform $O(\delta)$ trajectory bound, and contraction
improves the remainder to $O(\delta^2)$.  Applying the conclusion
with $\delta=|d|$ gives the stated pointwise bounds.
The derivative at $d=0$ is exactly the unchanged gradient
singleton, so the uniform analytic bounds give
$T_j^g=d+O(d^2)$, with a uniform derivative bound for the error.
Taking $d_0$ smaller makes $|(T_j^g)'(d)-1|<1/2$.
The quantitative analytic inverse theorem gives the common
inverse disk.  Exactness of the Gaussian integrations shows
that the trajectory with initial parameter $d$ maps to the
trajectory with that same parameter at the next step.
This proves the last assertion.  Finally
$(\mathcal R_j^g)'(0)=0$, so the contracted linear remainder
map is unchanged by \eqref{ep:normal-coordinates}.
\end{proof}

\begin{lemma}[Finite-volume response inequality]
\label{ep:finite-response}
For a pinned periodic height law with precision $A_0/k$,
let $C$ denote its covariance.  For every real $f_0$
annihilating constants,
\begin{equation}\label{ep:response-bound}
 \frac{\partial}{\partial k}\Var(h,f_0)
 \ge\frac{(Cf_0,A_0Cf_0)}{k^2}
 \ge\frac{\Var(h,f_0)^2}{k^2(f_0,A_0^+f_0)}.
\end{equation}
In particular, if $U_n=\Var(X_n)/g$ at an input with
$a(k)>0$, then
\begin{equation}\label{ep:positive-physical-response}
 \liminf_{n\to\infty}U_n'(k)\ge a(k)^2/k^2>0.
\end{equation}
The same strictly positive lower bound, weakened by a fixed
factor, holds uniformly on a sufficiently small neighborhood
where the limiting variance is bounded below.
\end{lemma}

\begin{proof}
Differentiate the absolutely convergent pinned lattice sum:
\[
 \partial_k\Var(h,f_0)
  =\frac1{2k^2}\Cov\bigl((h,f_0)^2,(h,A_0h)\bigr).
\]
Express the positive form $A_0$ as a sum of squares of real
linear forms $\ell_r$.  The paired fourth-cumulant inequality
in Lemma~\ref{st:finite-gaussian} gives
\[
 \Cov\bigl((h,f_0)^2,\ell_r(h)^2\bigr)
 \ge 2\Cov\bigl((h,f_0),\ell_r(h)\bigr)^2.
\]
Summing gives the first inequality in
\eqref{ep:response-bound}.  On the quotient by constants,
Cauchy--Schwarz gives
$(f_0,Cf_0)^2\le(f_0,A_0^+f_0)(Cf_0,A_0Cf_0)$;
pinning selects a representative and changes neither side.

Take $f_0=B_s^*f_n$.  Its bare Green form tends to $g$.
The observation-noise variance is independent of $k$ and
tends to zero, while the height variance tends to $a(k)g$.
Division by $g$ proves
\eqref{ep:positive-physical-response}.  The neighborhood
version follows from a positive lower variance bound and
continuity of $k$ and of the fixed bare denominator.
\end{proof}

\begin{proposition}[Openness and differentiability above $8\pi$]
\label{ep:supermarginal}
If $a(k_h)>8\pi$, then $k_h$ has an open neighborhood on
which $a>8\pi$ and $a$ is continuously differentiable.
Consequently
\begin{equation}\label{ep:endpoint-value}
 a(k_c)=8\pi .
\end{equation}
\end{proposition}

\begin{proof}
Use the gapped scheme without $Z$ and the coordinates of
Lemma~\ref{ep:gaussian-curves}.  Analyticity, exact vanishing
on $B=0$, and the diagonal balanced linearization give
\begin{align}
 \|B_+\|&\le[\theta+C(|u|+\|B\|)]\|B\|,\label{ep:gapped-B}\\
 |u_+-u|&\le C(|u|+\|B\|)\|B\|.\label{ep:gapped-u}
\end{align}
On a smaller tube the first factor is at most $\theta'<1$.
These inequalities keep nearby entry data in the tube,
make $B$ decay geometrically, and give a uniformly convergent
limit $u_\infty(k,a')$.  For fixed $a'$ it is analytic in $k$
on a complex neighborhood.  The convergence of its $k$
derivatives follows by Cauchy's formula on a smaller disk.

The scalar limit is jointly continuous in the real parameters.
To verify the point concerning changing frequency, truncate the
sole gradient balance sum in Equation~\eqref{wh:T-balance} after
$M$ terms. Its remaining operator norm is $O(\theta^M)$,
uniformly on the chosen compact reference interval in
$(8\pi,\infty)$; this scheme has no $Z$ balance sum.
Every retained scalar evaluation is a
finite-support integral or a finite Taylor projection.
Absolute expansions and Gaussian domination with reserve give
its continuity under varying $a'$.  First take the parameter
limit at fixed $M$, then $M\to\infty$.  The tail of
\eqref{ep:gapped-u} is uniformly summable.  The same argument,
with Cauchy's formula in $k$, proves continuity of
$\partial_k u_\infty$.  No norm continuity of periodic
activities as their period varies is used.

At $(k_h,a_h)$, Proposition~\ref{ep:no-exit} gives
$u_\infty=0$.  Terminal differentiation at a zero gives
\begin{equation}\label{ep:gapped-response}
 U_{n_j}'(k)\longrightarrow a'\partial_k u_\infty(k,a').
\end{equation}
Here $U_n$ is the physical variance defined in
Lemma~\ref{ep:finite-response}, with the same $s$, $f_n$, and
sequence $n_j$ throughout the neighborhood.  To check
\eqref{ep:gapped-response}, undo the balance and Gaussian
curve changes: the gradient singleton derivative tends to
$\partial_k u_\infty$, and the remainder derivatives tend
to zero by differentiated contraction.  The torus remainder
has its own differentiated contraction with these same
singleton inputs and also tends to zero.  Differentiating
\eqref{ep:gradient-diagnostic} twice in $w$ and once in $k$
leaves $a'\partial_k u_\infty$ by
\eqref{ep:F-energy}.  All discarded terms converge uniformly
on compact subsets of the zero set.  This argument computes
the derivative of the actual finite-volume observable and
does not differentiate a reference-dependent substitute.
Equation~\eqref{ep:positive-physical-response} implies
$\partial_k u_\infty(k_h,a_h)>0$.

Choose $k_-<k_h<k_+$ so close that $u_\infty(k_-,a_h)<0$
and $u_\infty(k_+,a_h)>0$, with
$\partial_k u_\infty>c>0$ on the intervening neighborhood.
Joint continuity preserves this bracket for every reference
$a'$ in an open interval about $a_h$.  The Intermediate Value Theorem supplies a unique zero $k=k(a')$ there.
At this zero the terminal activities vanish and the physical
variance tends to $a'g$; the structural limit identifies
$a(k(a'))=a'$.  Structural strict increase makes this
assignment unique even outside the chosen bracket.
Uniqueness and compactness show that $k(a')$ is continuous;
it is strictly increasing, and hence its image contains an
open interval about $k_h$.

Use one fixed entry construction and one fixed sequence of
physical functions $U_n$ on a compact subinterval of that
image.  At each input evaluate its derivative using the
appropriate reference $a'=a(k)$, held fixed during that
evaluation.  Uniform convergence in
\eqref{ep:gapped-response}, continuity of the zero assignment,
and continuity of its right-hand side imply uniform
convergence of $U_n'$ to a continuous function.  The
variances themselves converge to $a(k)$ uniformly there:
each $U_n$ is nondecreasing by Lemma~\ref{ep:finite-response},
and pointwise convergence to the continuous $a$ becomes uniform
by squeezing between the values on a finite grid whose successive
$a$ values differ by an arbitrarily small amount.
The Fundamental Theorem of Calculus, applied before the
limit and then passed under the uniform bound, proves that
$a$ is continuously differentiable, with derivative
$a(k)\partial_k u_\infty(k,a(k))$.

If $a(k_c)>8\pi$, the zero interval just obtained would
contain inputs below $k_c$ with positive coefficient.
This contradicts Corollary~\ref{an:physical-criterion}.
The positive-coefficient gap therefore forces
\eqref{ep:endpoint-value}.  Strict increase gives
$a(k)>8\pi$ for every $k>k_c$, so the differentiability
conclusion applies throughout that interval.
\end{proof}

\subsection{Critical state estimates}

The endpoint value and differentiability above it are now established.
To prove the infinite slope, we first need the decay rate of the actual
critical trajectory.  The free-face estimate from
Section~\ref{st:section} will rule out exponential decay of its odd
part; the local map will then determine its slower decay.

Initialize the two-coordinate maps at the actual onset, with the
fixed entry choices of Proposition~\ref{ep:no-exit}, and allow
true references $a'=a(k)$ for $k\ge k_c$ sufficiently close.
The entry-family balanced coordinates of
Equations~\eqref{wh:T-balance}--\eqref{wh:Z-balance} separate the
two scalar directions from the contracted remainder.
Equation~\eqref{ep:normal-coordinates} then parametrizes the exact
Gaussian curves by $u$ and measures the remainder relative to
those curves by $B$, leaving the fundamental coordinate $Z$
unchanged.  A common translation by
$\pi/\alpha$ splits the remainder as $B=B_e+B_o$ into its
even and odd parts.  This convention can be reversed once,
changing the sign of the fundamental coordinate, without
affecting the physical field.

\begin{lemma}[The parity-resolved map]
\label{ep:parity-map}
There are $\theta_1<1$, a fixed small tube, and a constant $C$
such that, uniformly in that tube and for the stated references,
\begin{align}
 u_+-u={}&-H_jZ^2+
 O\bigl((|u|+Z^2)(Z^2+\|B_e\|)+\|B_e\|^2
              +\|B_o\|(|Z|+\|B_o\|)\bigr),\label{ep:u-map}\\
 Z_+-\lambda Z={}&-b_juZ+
 O\bigl((u^2+Z^2+\|B_e\|)|Z|
 +( |u|+Z^2+\|B_e\|+\|B_o\|^2)\|B_o\|\bigr),\label{ep:Z-map}\\
 \|B_{e,+}\|\le{}&\theta_1\|B_e\|
                      +C(|Z|+\|B_o\|)^2,\label{ep:Be-map}\\
 \|B_{o,+}\|\le{}&\theta_1\|B_o\|
                      +C(|u|+Z^2+\|B_e\|)|Z|.\label{ep:Bo-map}
\end{align}
The remainders have the differentiated Taylor bounds implied
by these orders.  Here $H_j=-b_{j,{\rm en}}^{zz}$ and
$b_j=-b_{j,{\rm en}}^{tz}$ are the coefficients of the actual
entry-family balanced map.  They are uniformly bounded above and
bounded below by positive constants.  At the onset,
\[
 H_j\longrightarrow H:=H_{\rm en}
   =\frac{8\pi\cdot\pi}{4}e^{8\pi c_{v/\sqrt{32}}}\log L>0,
 \qquad
 b_j\longrightarrow b:=b_{\rm en}=2\log L>0 .
\]
All occurrences of $H,b$ and bounds for $H_j,b_j$ in the remaining
endpoint argument refer to this entry family.
\end{lemma}

\begin{proof}
The linear balanced map is diagonal on $(T,Z)$ and contracted
on the remainder.  Lemma~\ref{ep:gaussian-curves} removes
every term that survives on $Z=B=0$.
Half-period symmetry makes the first coordinate even and
the second odd.  Their surviving quadratic terms are
$-H_jZ^2$ and $-b_juZ$, with the entry-family signs and limits in
Proposition~\ref{wh:quadratic-coefficients}.  Indeed the change is
$T=T_j^g(u)=u+O(u^2)$ and
$\mathcal R=B+\mathcal R_j^g(u)=B+O(u^2)$, with both corrections
depending only on $u$ and with $Z$ unchanged.  The balanced linear
scalar map has no remainder forcing.  Substitution to quadratic order
can therefore change only the pure $u^2$ coefficient in the first
scalar equation; it preserves the coefficients of $Z^2$ and $uZ$.
The even remainder is forced at lowest order by two odd
factors.  The odd remainder has no linear fundamental forcing
by the diagonalization and is forced by $uZ$ and higher
orders.  Terms linear in a remainder, with a small
coefficient, are absorbed by increasing the contraction
constant from $\theta$ to $\theta_1<1$; small quadratic
terms involving $B_o$ in its own equation are absorbed there
as well.  Taylor expansion on a smaller analytic ball now
gives exactly \eqref{ep:u-map}--\eqref{ep:Bo-map}.
Cauchy's formula gives the corresponding differentiated
bounds.  The estimates remain valid with the entry
regulator by Lemma~\ref{en:regulator}.
\end{proof}

\begin{lemma}[The odd sector is not exponentially small]
\label{ep:odd-nonexponential}
At $k=k_c$, the pair $(Z_j,B_{o,j})$ does not satisfy
$|Z_j|+\|B_{o,j}\|\le C\vartheta^j$ for any
$\vartheta<1$.
\end{lemma}

\begin{proof}
Let $z_n^{\rm face}$ be the positive characteristic ratio
of the averaged four-face probability profile on a free
square of side $n$.  The critical coefficient is $8\pi$,
so Lemma~\ref{st:critical-face-ratio} gives a common
$c_0>0$ such that, for every fixed sufficiently large
$G=L^r$,
\[
 \liminf_{n\to\infty}
       z_{Gn}^{\rm face}/z_n^{\rm face}\ge c_0.
\]
Replace $c_0$ by a smaller number in $(0,1)$.
For each fixed $r$, the inequality with $c_0/2$ holds beyond
a threshold that may depend on $r$.  Iterating separately
on the $r$ residue classes of the index $j$ in
\eqref{ep:stop-sizes} gives
\[
 \limsup_{j\to\infty}
  \frac{-\log z_{n_j}^{\rm face}}j
       \le\frac{\log(2/c_0)}r .
\]
Since $r$ is arbitrary and $z_n^{\rm face}\le1$,
$-\log z_{n_j}^{\rm face}=o(j)$.

We spell out the change from a face profile to a uniform
profile.  Their difference has uniformly bounded bare
energy.  Along that direction write $\ell(t)=\log w(\rho+tq)$.
The logarithmic Hessian comparison of
Lemma~\ref{st:finite-gaussian} gives
$\ell''\ge-M$ for a fixed $M$, and $\ell\le0$.
Thus
$0\ge\ell(0)+t\ell'(0)-Mt^2/2$ for every real $t$,
so $|\ell'(0)|\le\sqrt{2M|\ell(0)|}$.
Integrating the Hessian bound from $0$ to $1$ gives
\[
 \ell(1)\ge\ell(0)-\sqrt{2M|\ell(0)|}-M/2 .
\]
It follows that the uniform-profile characteristic also
has negative logarithm $o(j)$.  Restoring the periodic bonds
increases this characteristic ratio by precision comparison.
The independent mean observation noise contributes the
positive factor $e^{-\Var(\bar\eta)/2}$, whose logarithm
is bounded.  Therefore
\begin{equation}\label{ep:phase-lower}
 -\log\E e^{i\bar Y_{n_j}}=o(j).
\end{equation}

If the plane odd sector decayed exponentially, the linear
balance relation and the contracted torus odd recursion
would make its unbalanced fundamental singleton and stopped
torus odd remainder exponentially small as well.  To see
the latter without identifying the two remainders, iterate
its contraction: it is bounded by an exponentially decaying
initial term plus a convolution of a fixed geometric kernel
with the exponentially small plane odd forcing.
In the terminal integral for $e^{i\bar Y}$, every nonzero
term contains an odd factor, since the common-mean integral
otherwise vanishes.  Absolute terminal bounds on the
remaining factors would then give
$|\E e^{i\bar Y_{n_j}}|\le C\vartheta_1^j$
for some $\vartheta_1<1$, contrary to
\eqref{ep:phase-lower}.
\end{proof}

\begin{lemma}[The compulsory state strip]
\label{ep:state-strip}
At $k_c$, after a finite step and a choice of the
half-period convention,
\[
 Z_j>0,\quad \|B_{o,j}\|\le\eps Z_j,\quad
 \|B_{e,j}\|\le C_eZ_j^2 .
\]
The first two remainder inequalities persist in a sufficiently
small tube, also for sufficiently small
$\mu=1-\lambda\ge0$.  Any trajectory in these cones
which tends to zero satisfies
\begin{equation}\label{ep:strip}
 c_-\frac{Z_j^2}{\mu+Z_j}
     <u_j<
 c_+\frac{Z_j^2}{\mu+Z_j},
 \qquad \|B_j\|=O(Z_j^2)
\end{equation}
after a further transient, with fixed $0<c_-<c_+$.
At $k_c$ in particular,
\begin{equation}\label{ep:critical-state-asymptotics}
 Z_j\asymp j^{-1},\qquad
 \sum_jZ_j=\infty,\qquad
 \frac{u_j}{Z_j}\longrightarrow\sqrt{H/b}.
\end{equation}
The strip in \eqref{ep:strip} is a necessary condition
for convergence to zero; its boundaries are not asserted
to be forward invariant.
\end{lemma}

\begin{proof}
Fix $\eps>0$ small.  When
$\|B_o\|\le\eps|Z|$, Equation~\eqref{ep:Z-map}
gives $|Z_+|\ge(1-o(1))|Z|$ at criticality, with preserved
sign, whereas \eqref{ep:Bo-map} gives
$\|B_{o,+}\|\le\theta_1\eps|Z|+o(|Z|)$.
Thus this cone is invariant in a small tube, and its ratio
tends to zero.  If the cone is never reached on a tail,
$|Z|<\eps^{-1}\|B_o\|$ there; Equation~\eqref{ep:Bo-map} gives
$\|B_{o,+}\|\le(\theta_1+C\delta/\eps)\|B_o\|$
when the tube has radius $\delta$.  Choosing
$\theta_1+C\delta/\eps<1$, and using
$|Z|<\eps^{-1}\|B_o\|$ at each later step of this
alternative, contracts the whole odd pair geometrically.
This contradicts Lemma~\ref{ep:odd-nonexponential}.
The odd pair cannot become zero, for then it would stay zero.
Reverse the convention if necessary to make $Z$ positive.

Its step ratio now tends to one.  Iterating
\eqref{ep:Be-map} against this slowly varying sequence gives
$\|B_e\|\le C_eZ^2$: choose a number $r<1$ with
$r^2>\theta_1$, then eventually $Z_+\ge rZ$,
and sum the geometric convolution of $Z_i^2$.
The homogeneous term is negligible since
$\log Z_j=o(j)$.  Taking $C_e$ larger makes this even
cone invariant whenever the step ratio remains close to
one; this applies uniformly for small $\mu$ and a small tube.

Within these cones, making $\eps$ and the tube smaller,
\eqref{ep:u-map}--\eqref{ep:Z-map} have the form
\begin{equation}\label{ep:reduced-state}
 u_+-u=-\widetilde H_j Z^2,\qquad
 Z_+/Z=1-\mu-\widetilde b_j u+O(Z^2),
\end{equation}
where $0<h_-\le\widetilde H_j\le h_+$ and
$0<b_-\le\widetilde b_j\le b_+$.
A trajectory converging to zero therefore has
$u_j=\sum_{i\ge j}\widetilde H_iZ_i^2>0$.

For $c>0$ put $f_c(Z)=cZ^2/(\mu+Z)$.
For a step with $u=f_c(Z)$, the Mean Value Theorem,
\eqref{ep:reduced-state}, and
$f_c'(Z)=cZ(2\mu+Z)/(\mu+Z)^2$ show that
\begin{equation}\label{ep:barrier-decrement}
 \frac{f_c(Z)-f_c(Z_+)}{Z^2}
 =
 c\frac{2\mu+Z}{\mu+Z}
 \left[\frac{\mu}{\mu+Z}
       +\widetilde b_jc\frac{Z^2}{(\mu+Z)^2}\right]
 +o(1),
\end{equation}
uniformly for fixed $c$ in the small tube.
The error is uniform also when $\mu/Z$ tends to zero or
infinity.  The prefactor $(2\mu+Z)/(\mu+Z)$ lies in
$[1,2]$.  If $c=c_-$ is sufficiently small, the right side
is less than $h_-/2$; if $c=c_+$ is sufficiently large, it
exceeds $2h_+$.  For the latter assertion write
$r=\mu/(\mu+Z)$: the bracket is
$r+\widetilde b_jc(1-r)^2$, whose minimum yields a
positive lower bound after multiplication by large $c$.
The choice of tube is made after fixing $c_\pm$.

The gap estimates on the exterior regions use the uniform
coefficient bounds, rather than monotonicity of the effective
coefficients. Write
$Z_+=Z(1-q)$, where
$q=\mu+\widetilde b_j u+e_j$ and $|e_j|\le C_0Z^2$,
and shrink the tube so $|q|\le1/2$.
If $0<u\le f_{c_-}(Z)$, then for $q<0$ the barrier
decrement is negative, whereas for $q\ge0$ convexity gives
\[
 \frac{f_{c_-}(Z)-f_{c_-}(Z_+)}{Z^2}
 \le 2c_-(1+b_+c_-)+2C_0c_-Z\le h_-/2 .
\]
Here choose $c_-$ so that the first term is at most
$h_-/4$, and then the tube so the second is at most
$h_-/4$. Thus the lower gap decreases by at least
$h_-Z^2/2$ throughout that exterior region.

For the upper region choose $c_+$ so large that
$\min\{c_+/4,b_-c_+^2/16\}\ge2h_+$, and then the
tube so $C_0(\mu+Z)\le b_-c_+/2$.
If $u\ge f_{c_+}(Z)$, then
$q\ge\mu+(b_-/2)u>0$.
The identity
$f_c'(y)/y=c(2\mu+y)/(\mu+y)^2$ shows that
$f_c'(y)\ge(y/Z)f_c'(Z)$ for $0<y\le Z$.
Since $Z_+\ge Z/2$, the Mean Value Theorem gives,
with $r=\mu/(\mu+Z)$,
\[
 \frac{f_{c_+}(Z)-f_{c_+}(Z_+)}{Z^2}
 \ge\frac{c_+}{2}
      \left[r+\frac{b_-c_+}{2}(1-r)^2\right]
 \ge\min\{c_+/4,b_-c_+^2/16\}\ge2h_+ .
\]
Hence the upper gap increases by at least $h_+Z^2$
throughout its exterior region.

A trajectory converging to zero has $u>0$ at every
such step. Starting on either barrier therefore gives
a strictly signed gap at the next step, and the
corresponding exterior gap inequality keeps that gap
bounded away from zero afterwards. This contradicts
$u,Z\to0$. The strict inequalities in
Equation~\eqref{ep:strip} follow. This argument does
not assert forward invariance of the strip.

In particular $u\le C Z$.  Equation~\eqref{ep:Bo-map}
now has forcing $O(Z^2)$, and its geometric convolution,
using the step ratio close to one, gives
$\|B_o\|=O(Z^2)$ after a transient.
At $\mu=0$ the strip gives $u\asymp Z$, so
\eqref{ep:reduced-state} gives
$cZ\le1-Z_+/Z\le CZ$.  Hence
$1/Z_+-1/Z$ is between two positive constants eventually,
which proves $Z_j\asymp j^{-1}$ and divergence of its sum.
Since the remainders are $O(Z^2)$ and the coefficients
converge, $x_j=u_j/Z_j$ satisfies
\[
 x_{j+1}-x_j=(bx_j^2-H+o(1))Z_j .
\]
It lies between two positive constants by the strip.
For any $\delta>0$, once the error is small, at or above
$\sqrt{H/b}+\delta$ its increment is at least $c_\delta Z_j$
and cannot cross that level downwards.  Divergence of
$\sum Z_j$ would force violation of the upper bound.
Below $\sqrt{H/b}-\delta$ the reverse argument violates
the positive lower bound.  Consequently $x_j$ tends to
$\sqrt{H/b}$.
\end{proof}

\subsection{Critical derivative growth}

The critical state estimate gives $Z_j\asymp j^{-1}$, so its cumulative
effect over scales is not summable.  We now differentiate the finite
maps to determine how this affects the response to the physical input.
The finite-volume response inequality will select the growing
derivative direction and its sign.

Let
\[
 p_j=\partial_k u_j,\qquad d_j=\partial_k Z_j,\qquad
 r_j=\partial_k B_j
\]
with $a'$ held fixed at each input.  These derivatives exist
at each finite step by Theorem~\ref{en:entry} and the
analytic maps.  They are not derivatives along the curve
$k\mapsto(k,a(k))$.

\begin{lemma}[Stable or growing derivative]
\label{ep:derivative-dichotomy}
At the critical input the derivative either tends to zero
in all coordinates, or its scalar part aligns with the
unstable eigenspace of
\begin{equation}\label{ep:critical-matrix}
 M=\begin{pmatrix}0&-2H\\-b&-\sqrt{bH}\end{pmatrix}.
\end{equation}
The second alternative actually occurs.  More precisely,
eventually
\begin{equation}\label{ep:critical-response-sign}
 p_j>0,\qquad p_j\longrightarrow\infty,\qquad
 \frac{d_j}{p_j}\longrightarrow-\frac12\sqrt{b/H},
 \qquad \|r_j\|=o(p_j).
\end{equation}
\end{lemma}

\begin{proof}
Differentiate the map in Lemma~\ref{ep:parity-map}.
Equation~\eqref{ep:critical-state-asymptotics} and
$\|B\|=O(Z^2)$ give
\begin{align}
 \binom{p_+}{d_+}
 &=\bigl[I+Z(M+o(1))\bigr]\binom p d
       +O(Z\|r\|),\label{ep:critical-derivative}\\
 \|r_+\|&\le\theta_2\|r\|+
                     CZ(|p|+|d|),\qquad \theta_2<1.
                     \label{ep:derivative-remainder}
\end{align}
In particular there is no derivative of $\lambda$ here.

Put $d_0=\sqrt{bH}$.  The eigenvalues of $M$ are
$d_0$ and $-2d_0$, with an unstable eigenvector
$(1,-\frac12\sqrt{b/H})$.  In scalar eigenvector
coordinates write $\xi$ for the unstable component and
$w$ for the stable one.  Choose $C_1$ large and set
$y=\max(|w|,C_1\|r\|)$.
For any prescribed small $\eps>0$, the cone
$y\le\eps|\xi|$ is invariant on a sufficiently late tail.
Indeed the stable scalar component has multiplier
$1-2d_0Z+o(Z)$, the unstable component has multiplier
$1+d_0Z+o(Z)$, and the $r$ contribution to either scalar
is at most $CZy/C_1$.  Choose $C_1$ so this coefficient
is smaller than $d_0/10$, then take the tail perturbation
smaller than a fixed multiple of $\eps d_0$.
The $r$ component contracts by $\theta_2$ with forcing
$O(Z|\xi|)$; taking the tail still later makes it satisfy
the same strict cone inequality.  Inside the cone
$|\xi_+|\ge(1+d_0Z/2)|\xi|$.

If the cone is never entered on this tail,
$|\xi|\le y/\eps$.  The stable estimate, after the above
choice of $C_1$ and a sufficiently small tail perturbation,
and the fast remainder estimate give
$y_+\le(1-cZ)y$ as long as this alternative persists.
For the remainder part use
$\theta_2y+CZy/\eps\le(1-cZ)y$ at sufficiently small $Z$.
Thus $y\to0$ and $\xi\to0$, since $\sum Z=\infty$.
If the cone is entered, it remains entered and $|\xi|$
diverges.  Applying the same comparison with arbitrarily
small prescribed $\eps$ shows $y/|\xi|\to0$.
This proves the dichotomy and alignment.

We now use the physical response to select its alternative
and sign.  Undoing the coordinate changes gives
$t_j'=p_j+O(u_jp_j)+O(\|r_j\|)$; analogous bounds hold
for the fundamental derivative.  A stopped torus remainder
derivative satisfies, separately from the plane remainder,
\begin{equation}\label{ep:torus-derivative-convolution}
 \|\partial_k\mathcal R_j^{\rm tor}\|
 \le C\theta_2^{j-j_0}
 +C\sum_{i=j_0}^{j-1}\theta_2^{j-1-i}
       \delta_i(|t_i'|+|z_i'|),
 \qquad \delta_i\longrightarrow0 .
\end{equation}
The update of the unbalanced remainder has no linear forcing
from a singleton.  Its differentiated forcing therefore has
the coefficient
$O(|t_i|+|z_i|+\|\mathcal R_i^{\rm tor}\|)$,
which tends to zero by the state contraction.  This proves
\eqref{ep:torus-derivative-convolution}, with the bounded
initial segment absorbed into its first term.

If all derivatives tended to zero, the convolution would
tend to zero and terminal differentiation would give
$U_{n_j}'(k_c)\to0$, contradicting
\eqref{ep:positive-physical-response}.  In the growing
alternative the cone bounds the magnitude of every earlier
scalar derivative on the growing tail by a fixed multiple
of its current magnitude.  Split the convolution into a
fixed initial segment and a tail on which $\delta_i$ is
arbitrarily small.  It follows that its right side is
$o(|p_j|)$.  The same holds for the plane remainder.
Differentiating the terminal logarithmic ratio twice in the
source and once in $k$ consequently gives
\begin{equation}\label{ep:critical-terminal-response}
 U_{n_j}'(k_c)=a(k_c)p_j+o(|p_j|).
\end{equation}
Linear odd terms vanish in this mean-free diagnostic.
The positive lower bound in
\eqref{ep:positive-physical-response} forces the sign in
\eqref{ep:critical-terminal-response} to be positive.
The unstable eigenvector then gives exactly
\eqref{ep:critical-response-sign}.
\end{proof}

\begin{lemma}[Entering the cones at nearby references]
\label{ep:parameter-continuity}
Fix $M_*$ large and $\eta_*>0$ small with strict inequalities
$\eta_*<\tfrac12\sqrt{b/H}<M_*$, and fix any $\eta_1>0$.
A single sufficiently late finite step can be chosen so
that, for every sufficiently close warmer true reference,
the state enters the cones of Lemma~\ref{ep:state-strip}
and the derivative satisfies
\[
 p>0,\qquad -M_*<d/p<-\eta_*<0,\qquad
 \|r\|<\eta_1p .
\]
The scalar inequalities in \eqref{ep:strip} and a bound
$\|B\|\le C_BZ^2$ hold from this step onward, with $C_B$ depending
only on the coarse state estimates, independently of the chosen step,
prescribed tube, and warmer reference.
The step can also be required to lie in any prescribed sufficiently
small state tube.  This conclusion uses no continuity in the full
activity norm as $a'$ varies.
\end{lemma}

\begin{proof}
For every fixed step the scalar coordinates and their
held-reference $k$ derivatives converge to their critical
values as $k\downarrow k_c$.  Here is the precise
continuity argument.  The initial finite-support integrals
are continuous in $(k,a')$: at the fixed entry mesh the
massive energy matrices have uniform coercivity, and all
link sums have uniformly summable truncations.  A finite
number of convolutions and local projections can be tested
on fixed fields and directions, including the constant,
affine, and Gaussian arguments of the localization rules.
Gaussian domination with the reserve in the $W e^V$
regulator passes the limit in each such integral.
For any fixed depth there is a locally uniform Cauchy disk
in $k$, so the same statement holds for its input
derivatives.  The radius is allowed to depend on that
finite depth.  Truncate the two balance sums in
Equations~\eqref{wh:T-balance} and~\eqref{wh:Z-balance}.
Here the reference interval is chosen with
$\lambda\ge\lambda_*:=1-\mu_0>\theta$, so their remaining
operator norms are respectively $O(\theta^M)$ and
$O((\theta/\lambda_*)^M)$, uniformly on that interval.
Remove these tails after taking the parameter limit at
fixed $M$.
The analytic inverse
Gaussian coordinates preserve scalar continuity.
These operations prove the asserted finite-step scalar
continuity without comparing functions of different
periods in a norm.

For the remainder norms use upper limits at each fixed
step.  They are initially bounded, including their input
derivatives, by the uniform entry estimates.  Their upper
limits obey \eqref{ep:Be-map}--\eqref{ep:Bo-map} along
the limiting critical scalar trajectory.
Contraction first gives an odd bound $O(Z)$; substituting
that bound gives an even bound $O(Z^2)$, and substituting
$u=O(Z)$ and that even bound gives the odd bound $O(Z^2)$.
The homogeneous terms vanish relative to these powers
because the critical scalar sequence varies only on the
$1/j$ scale.  The coarse choice of $C_e$ can be made with a strict
margin over the eventual even upper-limit ratio, since both bounds
use the same geometric contraction.  Retain that fixed choice.

For the derivative upper limits,
\eqref{ep:derivative-remainder} gives a geometric
convolution forced by $CZ(|p|+|d|)$ along the critical
trajectory.  Its ratio to the growing critical $p$ tends
to zero, by the same split into a fixed initial history and
a late history as in
\eqref{ep:torus-derivative-convolution}.

The quadratic remainder constant can also be fixed independently of
the selected step.  Let $q_0\in(\sqrt{\theta_1},1)$ be a fixed
lower bound for $Z_+/Z$ in the persistent remainder cones of the
coarse tube and reference range, as in the proof of
Lemma~\ref{ep:state-strip}.  Each nearby true-reference trajectory
stays in that tube and
tends to zero by Proposition~\ref{ep:no-exit}.  Applied to any tail
in the cones, the exterior-gap argument proving \eqref{ep:strip}
therefore gives $0<u<c_+Z$ at every step of that tail.  The even cone
and \eqref{ep:Bo-map} consequently give
\[
 \|B_{o,+}\|\le\theta_1\|B_o\|+FZ^2
\]
with $F$ fixed by the coarse state estimates.  The preceding
upper-limit convolution also gives a fixed eventual bound $D_o$ for
$\limsup_{k\downarrow k_c}\|B_{o,j}(k)\|/Z_j(k_c)^2$.
Choose $C_o$ larger than both $D_o$ and
$F/(q_0^2-\theta_1)$, with strict margins.  Then
$\theta_1C_o+F<q_0^2C_o$, so the displayed recursion preserves
$\|B_o\|\le C_oZ^2$.  Together with the even cone this gives the
fixed constant $C_B=C_e+C_o$.

Now choose a finite step so late that all state and derivative
inequalities, including the fixed quadratic remainder bound, have
strict margins there and the state lies in the prescribed small
tube.  Finite-step scalar continuity and the remainder upper-limit
bounds preserve those margins for $k>k_c$ sufficiently close.
The invariant remainder cones persist, the scalar strip applies
to their convergent tails, and the fixed quadratic bound just proved
persists at every later step.
\end{proof}

\subsection{The warmer crossover and the infinite endpoint slope}

The preceding growth concerns derivatives at the critical input.
To identify the slope of the limiting coefficient, we must also follow
nearby warmer trajectories through their eventual decay.  We first take
their terminal-volume limit at each fixed input, then let that input
decrease to $k_c$.  The crossover estimate below shows that the growth
accumulated before decay persists in this order of limits.

\begin{proposition}[Divergent response from the warmer side]
\label{ep:infinite-response}
As $k\downarrow k_c$,
\begin{equation}\label{ep:infinite-derivative}
 \frac{d}{dk}a(k)\longrightarrow+\infty ,
 \qquad
 \frac{a(k)-a(k_c)}{k-k_c}\longrightarrow+\infty .
\end{equation}
\end{proposition}

\begin{proof}
The fixed late entry step from Lemma~\ref{ep:parameter-continuity}
will be chosen after the constants below.  For a warmer reference
put $\mu=1-L^{2-a(k)/(4\pi)}>0$, held fixed in each
input differentiation.  Work first in the fixed coarse tube of the
preceding state argument.  For a convergent tail in its persistent
remainder cones, the reduced state equation and scalar strip give,
uniformly in this tube,
\begin{equation}\label{ep:crossover-state}
 c(\mu+Z)\le1-Z_+/Z\le C(\mu+Z).
\end{equation}
Indeed
\[
 \frac{\mu+Z^2/(\mu+Z)}{\mu+Z}
 =1-\frac{Z}{\mu+Z}+\left(\frac{Z}{\mu+Z}\right)^2
 \in[3/4,1].
\]
The $O(Z^2)$ error in
\eqref{ep:reduced-state} is absorbed by the earlier small choice of
the coarse tube.  That choice and the reference range also make the
right side of \eqref{ep:crossover-state} less than $1/2$, so $Z$
decreases along the whole tail.

Use the equivalent state size
$\sigma=|u|+|Z|+\|B_e\|+\|B_o\|$ for the tubes, and write the fixed
coarse tube as $\sigma<\rho_{\rm c}$.  The persistent remainder cones
and the upper scalar strip imply, at every step after cone entry,
\[
 \sigma_j\le
 (1+c_++\eps+C_e\rho_{\rm c})Z_j
 =:KZ_j\le KZ_{j_0}\le K\sigma_{j_0},\qquad j\ge j_0 .
\]
Here $K$ depends only on the coarse state estimates, not on the late
step or the smaller derivative tube chosen below.  This tail bound
uses no differentiated estimate.
At any chosen fixed late step, scalar continuity makes the initial
$Z$ bounded above and bounded below by positive constants uniformly
as $k\downarrow k_c$.

We record the derivative comparison with its parameter
order.  Fix positive bounds $H_\pm,b_\pm$ for the entry-family
coefficients.  Choose $M_*>\tfrac12\sqrt{b/H}$ large and then $\eta>0$ small so that
the two broad endpoint signs verified below have strict margins.
While $Z\ge\mu$, the state strip gives
$\mu+b_ju\le C_*Z$ for a fixed $C_*$.  Next choose $\eta_*>0$
so small that
\[
 \eta_*<\tfrac12\sqrt{b/H},\qquad
 2H_+\eta_*^2+C_*\eta_*<b_-/2 .
\]
Choose $\eps_1>0$ after these constants, small compared with all
the strict endpoint margins and with $H_-\eta_*$.  Then choose the
remainder ratio $\eta_1>0$ and the state tube sufficiently small
that the differentiated Taylor errors below have size
$O(\eps_1Z)$ and the strict remainder contraction preserves
$\|r\|\le\eta_1p$.  In particular, its radius
$\rho_{\rm d}\le\rho_{\rm c}$ in the state size $\sigma$ is chosen
after $\eta_1$, using the fixed $C_B$ from
Lemma~\ref{ep:parameter-continuity}.  Apply that lemma with the
smaller entry radius $\rho_{\rm d}/(2K)$.  It supplies a single late
entry step with $\sigma_{j_0}<\rho_{\rm d}/(2K)$, and the preceding
tail bound gives $\sigma_j<\rho_{\rm d}/2$ for every $j\ge j_0$.
Thus the differentiated estimates apply throughout the tail.  Shrink
the warmer neighborhood so that its initial $Z$ also exceeds
$\mu$.  Put $x=d/p$.  In the cone
\[
 -M_*\le x\le\eta,\qquad p>0,\qquad \|r\|\le\eta_1p ,
\]
the differentiated maps give
\begin{align}
 p_+/p&=1-2H_jZx+O(\eps_1Z),\label{ep:warm-p}\\
 d_+/p&=(1-\mu-b_ju)x-b_jZ+O(\eps_1Z),\label{ep:warm-d}
\end{align}
with the just chosen $\eps_1$.  Terms involving $B=O(Z^2)$ and higher scalar
orders are included in that error; those involving $r$
are bounded by $CZ\eta_1$.  Equation
\eqref{ep:derivative-remainder} remains valid.

The unperturbed ratio map in
\eqref{ep:warm-p}--\eqref{ep:warm-d} is increasing on this
fixed interval: its derivative has positive numerator
$1-\mu-b_ju-2H_jb_jZ^2$ and a positive squared denominator.
At an endpoint $x$ its increment has numerator
\begin{equation}\label{ep:ratio-numerator}
 2H_jZx^2-(\mu+b_ju)x-b_jZ .
\end{equation}
At $x=-M_*$ this is positive by at least $cZ$ if $M_*$
is large; at $x=\eta$ it is negative by at least $cZ$
if $\eta$ is small.  These strict signs absorb the errors
in \eqref{ep:warm-p}--\eqref{ep:warm-d}.
The denominator stays positive in the small tube.
The remainder inequality, its strict contraction, and
small $Z$ preserve $\|r\|\le\eta_1p$.
Thus the full derivative cone is invariant.

While $Z\ge\mu$, the upper boundary can be replaced by the already
chosen $x=-\eta_*$: Equation~\eqref{ep:ratio-numerator} there is
at most
$Z(2H_+\eta_*^2+C_*\eta_*-b_-)<-b_-Z/2$.
The critical entry value lies strictly below this
boundary by Lemma~\ref{ep:parameter-continuity}.
The choice of $\eps_1$ preserves this sign and makes
$2H_-\eta_*-O(\eps_1)>0$.  Consequently throughout this first stage,
\begin{equation}\label{ep:pre-crossover-growth}
 p_+\ge(1+cZ)p .
\end{equation}

Let $j_\times$ be the first index with $Z<\mu$.
It is finite by \eqref{ep:crossover-state}.
Before that index the relative decrement is comparable
to $Z$.  Taking logarithms in
\eqref{ep:crossover-state} therefore gives
\[
 \sum_{j_0\le j<j_\times}Z_j
 \ge c\log(Z_{j_0}/\mu)-C\longrightarrow\infty .
\]
The last step crosses $\mu$ by a ratio bounded below by
$1/2$, which accounts for the fixed additive constant.
The initial $p_{j_0}$ stays bounded below positively.
Equation~\eqref{ep:pre-crossover-growth} implies
$p_{j_\times}\to\infty$ as $\mu\downarrow0$.

After crossover, \eqref{ep:crossover-state} gives
$Z_{j+1}\le(1-c\mu)Z_j$, and hence
\begin{equation}\label{ep:post-crossover-sum}
 \sum_{j\ge j_\times}Z_j
 \le Z_{j_\times}/(c\mu)\le C .
\end{equation}
The larger invariant ratio cone gives
$p_+/p\ge1-CZ$, so the possible decrease of $p$
on this entire tail is at most a fixed multiplicative
factor.  For each fixed $\mu>0$,
$|p_+/p-1|\le CZ$ and summability imply the existence
of a finite positive limit $p_\infty$.
The equation for $d$ has a multiplier tending to
$1-\mu<1$ and a forcing tending to zero; the contracted
remainder equation has the same property.  Thus
\begin{equation}\label{ep:warm-limits}
 p_j\to p_\infty>0,\qquad d_j\to0,\qquad r_j\to0,
 \qquad p_\infty\to\infty\quad(\mu\downarrow0).
\end{equation}
For the last limit combine the divergent pre-crossover
gain with \eqref{ep:post-crossover-sum}.

On compact subintervals of $k>k_c$ sufficiently close
to $k_c$, the parameter $\mu$ is bounded away from zero
and is continuous by Proposition~\ref{ep:supermarginal}.
All tail estimates just used are then uniform.
Undoing the coordinates and applying the differentiated
terminal diagnostic, including the separate torus
convolution estimate, yields
\[
 U_{n_j}'(k)\longrightarrow a(k)p_\infty(k)
\]
uniformly on those compact subintervals.
This uses the physical functions $U_{n_j}$ of the current entry
construction.  Each is continuously differentiable and nondecreasing
in $k$ by Lemma~\ref{ep:finite-response}, and the structural limit
gives $U_{n_j}(k)\to a(k)$ pointwise.  On a compact warmer interval,
$a$ is continuous by Proposition~\ref{ep:supermarginal}.  A finite
grid on which its successive increments are arbitrarily small,
together with monotonicity of each $U_{n_j}$, therefore makes this
convergence uniform.  Write $v(k)=a(k)p_\infty(k)$.  The displayed
uniform derivative convergence makes $v$ continuous, as the uniform
limit of the finite continuous derivatives.  Applying the Fundamental
Theorem of Calculus to this same sequence and passing to the limit gives
\[
 a(y)-a(x)=\int_x^y v(k)\,\dd k
\]
for every $x<y$ in the compact interval.  Thus
$\frac{d}{dk}a(k)=v(k)$ there, and
\eqref{ep:warm-limits} proves the first limit in
\eqref{ep:infinite-derivative}.  Finally $a$ is right
continuous at $k_c$ and differentiable on its right.
For any $M>0$, the derivative exceeds $M$ on a sufficiently
short right interval.  Apply the Mean Value Theorem on
$[k_c,k]$ to obtain
$(a(k)-a(k_c))/(k-k_c)\ge M$ there.
Since $M$ is arbitrary, the second limit follows.
\end{proof}

\begin{corollary}[Completion of the height assertions]
\label{ep:height-conclusion}
For every admissible finite-range $J$, the uniquely determined
effective temperature is
\[
 \beff(J,\beta)=v_J^2a(\beta/v_J^2),\qquad
 \beta\ge\bc(J).
\]
It has the field, Laplace-functional, and tightness conclusions
in Theorem~\ref{main:heights}, is differentiable for $\beta>\bc(J)$, and satisfies
\[
 \beff(J,\bc(J))=8\pi v_J^2,\qquad
 \lim_{\beta\downarrow\bc(J)}
 \frac{\beff(J,\beta)-\beff(J,\bc(J))}
      {\beta-\bc(J)}=+\infty .
\]
Its high-temperature normalization and spread-out onset are
\[
 \beff(J,\beta)=\beta+O_J(e^{-c_J\beta}),\qquad
 \frac{\bc(J_\rho)}{8\pi v_{J_\rho}^2}\longrightarrow1 .
\]
For nearest neighbors the endpoint value is $2\pi$.
\end{corollary}

\begin{proof}
The physical criterion, including the endpoint, is
Corollary~\ref{an:physical-criterion}.  Its torus coefficient
and full Laplace and distributional limits are
Theorem~\ref{an:dirichlet-torus}; a single nonzero mean-free
test determines that coefficient uniquely.
Those structural limits hold along all growing square
sides, so they apply to every sequence $n=L^N$, with for
example $L=2$, independently of the larger auxiliary
factor used in the endpoint maps.
Equation~\eqref{ep:endpoint-value} and
Propositions~\ref{ep:supermarginal} and~\ref{ep:infinite-response}
give the endpoint and regularity statements.  The change
of variables $k=\beta/v_J^2$ cancels the two factors
$v_J^2$ in the derivative and in the difference quotient.
Corollaries~\ref{wh:high-temperature} and~\ref{wh:spread-out}
supply the last two asymptotics with exactly this
normalization.  Finally $v_{J_{\nn}}^2=1/4$, so
$8\pi v_{J_{\nn}}^2=2\pi$.
\end{proof}
